%%% Stoffe und ihre Eigenschaften
\begin{lesen}
\lies{\skElement \; \skVerbindung}{An element consists of only one kind of atom; a compound consists of two or more elements chemically combined.}{Ein Element besteht aus nur einer Atomsorte; eine Verbindung besteht aus zwei oder mehr chemisch verbundenen Elementen.}
\lies{\skLoesung}{Salt dissolves in water: salt is the solute, water is the solvent, and salt water is the solution.}{Salz löst sich in Wasser: Salz ist der gelöste Stoff, Wasser das Lösungsmittel und Salzwasser die Lösung.}
\lies{\skFiltration}{Sand can be separated from water by filtration: the sand stays in the filter paper as the residue, the water passes through as the filtrate.}{Sand lässt sich von Wasser durch Filtrieren trennen: Der Sand bleibt als Rückstand im Filter, das Wasser läuft als Filtrat durch.}
\lies{Eigenschaften}{Copper is a reddish-brown metal. It is a good conductor of electricity and heat and it melts at 1085 degrees Celsius.}{Kupfer ist ein rotbraunes Metall. Es leitet den elektrischen Strom und Wärme gut und schmilzt bei 1085 Grad Celsius.}
\lies{Destillation}{Ethanol and water can be separated by distillation because they have different boiling points.}{Ethanol und Wasser lassen sich durch Destillation trennen, weil sie verschiedene Siedetemperaturen haben.}
\end{lesen}

%%% Laborgeräte und Arbeitsweisen
\begin{lesen}
\lies{\skMesszylinder}{Measure out fifty millilitres of water using a measuring cylinder.}{Miss mit dem Messzylinder fünfzig Milliliter Wasser ab.}
\lies{\skErhitzen}{Heat the test tube gently in the Bunsen flame. Do not point the test tube at anyone.}{Erhitze das Reagenzglas vorsichtig in der Brennerflamme. Richte die Öffnung des Reagenzglases nicht auf Personen.}
\lies{\skBecherglas}{Pour the solution into a beaker and stir it with a glass rod.}{Gieße die Lösung in ein Becherglas und rühre mit einem Glasstab um.}
\lies{\skWaage}{Weigh out twelve point five grams of copper sulfate on the balance.}{Wiege auf der Waage zwölf Komma fünf Gramm Kupfersulfat ab.}
\lies{\skSchutzbrille}{Always wear safety goggles when you are working with chemicals.}{Trage beim Arbeiten mit Chemikalien immer eine Schutzbrille.}
\end{lesen}

%%% Atombau und Periodensystem
\begin{lesen}
\lies{\skSchalen}{A sodium atom has eleven protons in its nucleus and eleven electrons in three shells: two, eight, one.}{Ein Natriumatom hat elf Protonen im Kern und elf Elektronen auf drei Schalen: zwei, acht, eins.}
\lies{\ce{Na -> Na+ + e-}}{A sodium atom loses its outer electron and forms a positively charged sodium ion.}{Ein Natriumatom gibt sein Außenelektron ab und bildet ein positiv geladenes Natrium-Ion.}
\lies{\pse{(0,0) rectangle (0.3,1.2)}}{The alkali metals are in group one of the periodic table. Their atoms have one outer electron.}{Die Alkalimetalle stehen in der ersten Hauptgruppe des Periodensystems. Ihre Atome haben ein Außenelektron.}
\lies{\pse{(2.1,0) rectangle (2.4,1.5)}}{The noble gases have full outer shells and are very unreactive.}{Die Edelgase haben voll besetzte Außenschalen und sind sehr reaktionsträge.}
\lies{\ce{^{35}_{17}Cl}, \ce{^{37}_{17}Cl}}{Chlorine has two isotopes: chlorine-thirty-five and chlorine-thirty-seven. They have the same number of protons but different numbers of neutrons.}{Chlor hat zwei Isotope: Chlor-35 und Chlor-37. Sie haben gleich viele Protonen, aber verschieden viele Neutronen.}
\end{lesen}

%%% Chemische Bindung
\begin{lesen}
\lies{\skGitter}{In sodium chloride the sodium ions and the chloride ions are held together in a lattice by ionic bonds.}{Im Natriumchlorid werden die Natrium-Ionen und die Chlorid-Ionen durch Ionenbindung in einem Gitter zusammengehalten.}
\lies{\skAtombindung}{In a hydrogen molecule the two atoms share a pair of electrons: a covalent bond.}{Im Wasserstoffmolekül teilen sich die beiden Atome ein Elektronenpaar: eine Atombindung.}
\lies{$\mathrm{O{=}C{=}O}$}{A carbon dioxide molecule contains two double bonds.}{Ein Kohlenstoffdioxidmolekül enthält zwei Doppelbindungen.}
\lies{\skDipol}{The water molecule is a dipole: the oxygen atom carries a partial negative charge, the hydrogen atoms a partial positive charge.}{Das Wassermolekül ist ein Dipol: Das Sauerstoffatom trägt eine negative, die Wasserstoffatome tragen eine positive Teilladung.}
\lies{\skMetall}{Metals conduct electricity because their delocalised electrons are free to move.}{Metalle leiten den elektrischen Strom, weil ihre delokalisierten Elektronen frei beweglich sind.}
\end{lesen}

%%% Chemische Reaktionen
\begin{lesen}
\lies{\ce{2 Mg + O2 -> 2 MgO}}{Magnesium reacts with oxygen to form magnesium oxide. Magnesium and oxygen are the reactants, magnesium oxide is the product.}{Magnesium reagiert mit Sauerstoff zu Magnesiumoxid. Magnesium und Sauerstoff sind die Ausgangsstoffe, Magnesiumoxid ist das Reaktionsprodukt.}
\lies{\skExotherm}{The reaction is exothermic: energy is given out to the surroundings and the temperature rises.}{Die Reaktion ist exotherm: Es wird Energie an die Umgebung abgegeben und die Temperatur steigt.}
\lies{\skKatalysator}{A catalyst speeds up a reaction by lowering the activation energy. It is not used up in the reaction.}{Ein Katalysator beschleunigt eine Reaktion, indem er die Aktivierungsenergie herabsetzt. Er wird dabei nicht verbraucht.}
\lies{\ce{N2 + 3 H2 <=> 2 NH3}}{The reaction is reversible. At equilibrium the forward and the reverse reaction take place at the same rate.}{Die Reaktion ist umkehrbar. Im Gleichgewicht laufen Hin- und Rückreaktion gleich schnell ab.}
\lies{\skNiederschlag}{When silver nitrate solution is added, a white precipitate of silver chloride forms.}{Bei Zugabe von Silbernitratlösung bildet sich ein weißer Niederschlag von Silberchlorid.}
\end{lesen}

%%% Säuren, Basen und Salze
\begin{lesen}
\lies{\ce{HCl + NaOH -> NaCl + H2O}}{Hydrochloric acid reacts with sodium hydroxide solution to form sodium chloride and water: an acid and a base neutralise each other.}{Salzsäure reagiert mit Natronlauge zu Natriumchlorid und Wasser: Säure und Base neutralisieren einander.}
\lies{\skPH}{Acidic solutions have a pH below seven, alkaline solutions have a pH above seven.}{Saure Lösungen haben einen pH-Wert unter sieben, alkalische Lösungen einen pH-Wert über sieben.}
\lies{Lackmus}{Acids turn blue litmus red; alkalis turn red litmus blue.}{Säuren färben blaues Lackmus rot; Laugen färben rotes Lackmus blau.}
\lies{\skTitration}{The acid is titrated with sodium hydroxide solution until the indicator changes colour.}{Die Säure wird mit Natronlauge titriert, bis der Indikator umschlägt.}
\lies{\ce{CH3COOH}}{Ethanoic acid is a weak acid: it is only partially ionised in water.}{Essigsäure ist eine schwache Säure: Sie ist in Wasser nur teilweise dissoziiert.}
\end{lesen}

%%% Redoxreaktionen und Elektrochemie
\begin{lesen}
\lies{\ce{2 Cu + O2 -> 2 CuO}}{Copper is oxidised to copper(II) oxide: the copper atoms lose electrons.}{Kupfer wird zu Kupfer(II)-oxid oxidiert: Die Kupferatome geben Elektronen ab.}
\lies{\ce{CuO + H2 -> Cu + H2O}}{Copper(II) oxide is reduced to copper; hydrogen is the reducing agent.}{Kupfer(II)-oxid wird zu Kupfer reduziert; Wasserstoff ist das Reduktionsmittel.}
\lies{Merkhilfe}{OIL RIG: Oxidation Is Loss, Reduction Is Gain (of electrons).}{Oxidation ist Elektronenabgabe, Reduktion ist Elektronenaufnahme.}
\lies{\skElektrolyse}{During electrolysis the positive ions move to the cathode and the negative ions move to the anode.}{Bei der Elektrolyse wandern die positiven Ionen zur Kathode und die negativen Ionen zur Anode.}
\lies{\skRost}{Iron rusts in the presence of water and oxygen.}{Eisen rostet, wenn Wasser und Sauerstoff vorhanden sind.}
\end{lesen}

%%% Organische Chemie
\begin{lesen}
\lies{\ce{CH4 + 2 O2 -> CO2 + 2 H2O}}{Methane burns in oxygen to form carbon dioxide and water.}{Methan verbrennt mit Sauerstoff zu Kohlenstoffdioxid und Wasser.}
\lies{\ce{H2C=CH2}}{Ethene is an unsaturated hydrocarbon: its molecule contains a carbon–carbon double bond. It decolourises bromine water.}{Ethen ist ein ungesättigter Kohlenwasserstoff: Sein Molekül enthält eine C=C-Doppelbindung. Es entfärbt Bromwasser.}
\end{lesen}

%%% Stoffmenge und chemisches Rechnen
\begin{lesen}
\lies{$M(\mathrm{H_2O}) = 18\eh{g/mol}$}{The molar mass of water is eighteen grams per mole.}{Die molare Masse von Wasser beträgt achtzehn Gramm pro Mol.}
\lies{$n = \dfrac{m}{M} = \dfrac{36\eh{g}}{18\eh{g/mol}} = 2\eh{mol}$}{Thirty-six grams of water contain two moles of water molecules.}{Sechsunddreißig Gramm Wasser enthalten zwei Mol Wassermoleküle.}
\lies{$c = \dfrac{n}{V} = \dfrac{0{,}5\eh{mol}}{2\eh{l}} = 0{,}25\eh{mol/l}$}{The concentration of the solution is nought point two five moles per litre.}{Die Konzentration der Lösung beträgt null Komma zwei fünf Mol pro Liter.}
\lies{$V = n \cdot V_m = 2\eh{mol} \cdot 22{,}4\eh{l/mol}$ \newline $V = 44{,}8\eh{l}$}{Under standard conditions two moles of a gas occupy forty-four point eight litres.}{Unter Normbedingungen nehmen zwei Mol eines Gases vierundvierzig Komma acht Liter ein.}
\end{lesen}
